% =====================================================================================
% Figure -- the walkable region of the site, top-down.
% Belongs in \subsection{Object Layout Generation} (\label{sec:layouts}), after the
% sentence ending "...defines where an object may appear."
%
% No extra packages required.
%
% Expected image file, relative to the main .tex:
%     figures/walkable-area.jpg
%
% CAPTION CAVEAT -- read before using. methodology_section.tex (lines 281--284) records
% that the walkable region and the spawn region are deliberately SEPARATE objects in the
% scene: spawning uses RandomSpawner.targetPlanes, walk-only surfaces live on a
% WalkableAreaSet component. If the green region rendered here is the WalkableAreaSet and
% not targetPlanes, the caption must NOT call it the spawn domain -- use the alternate
% caption at the bottom of this file instead. Confirm which object produced the render.
% =====================================================================================

\begin{figure}[t]
  \centering
  \includegraphics[width=\linewidth]{figures/walkable-area}
  \caption{The walkable region of the study site (green), seen from above over the site
    reconstruction. The region was derived from the site scan by reducing the mesh to the
    traversable ground surface alone, removing buildings, vegetation, planting beds, and
    the central lawn, and totals 5854\,m\textsuperscript{2}. This trimmed geometry, rather
    than a bounding box or an authored polygon, defines where an object may appear:
    candidate positions falling outside it are discarded, so the mesh acts directly as the
    acceptance mask during layout generation.}
  \Description{A top-down orthographic view of a long rectangular campus courtyard,
    reconstructed from a photographic scan, with the walkable ground shaded in translucent
    green and outlined in dark green. The shaded region covers the paved walkways that run
    the length of the courtyard, a wide paved plaza at the left end surrounding a circular
    planting bed, and a paved area at the upper right. It excludes the buildings along both
    long edges, the large rectangular lawn in the middle of the courtyard, the planting beds
    and rows of shrubs and trees that separate the walkways, and the parking area with cars
    at the top. A legend in the lower right reads "Walkable area".}
  \label{fig:walkable}
\end{figure}

% -------------------------------------------------------------------------------------
% ALTERNATE CAPTION -- use this one if the green region is the WalkableAreaSet (the
% surface participants may traverse) and NOT RandomSpawner.targetPlanes (the surface
% objects spawn on). It describes the region without claiming it is the spawn domain.
% -------------------------------------------------------------------------------------
%
%   \caption{The walkable region of the study site (green), seen from above over the site
%     reconstruction. The region was derived from the site scan by reducing the mesh to the
%     traversable ground surface alone, removing buildings, vegetation, planting beds, and
%     the central lawn, and totals 5854\,m\textsuperscript{2}. It bounds where
%     participants may move, and supplies the surface over which shortest-path tours are
%     measured.}
